\documentclass[11pt,a4paper]{article}

\usepackage[utf8]{inputenc}
\usepackage[T1]{fontenc}
\usepackage[margin=1in]{geometry}
\usepackage{graphicx}
\usepackage{hyperref}
\usepackage{setspace}
\usepackage{enumitem}
\usepackage{titlesec}
\usepackage{float}

% Formatting 
\setstretch{1.15}
\titleformat{\section}{\large\bfseries}{\thesection.}{0.5em}{}
\titleformat{\subsection}{\normalsize\bfseries}{\thesubsection}{0.5em}{}
\setlength{\parindent}{0pt}
\setlength{\parskip}{0.8em}

\title{\vspace{-2cm}\textbf{Sortamine Project Report}\\ \large Data Structures and Algorithms}
\author{}
\date{\vspace{-1cm}}

\begin{document}

\maketitle
\vspace{-1cm}

\section{Design and Code Structure}

The Sortamine application is built using a modern \textbf{Model-View-Controller (MVC) architecture} implemented via \textbf{JavaFX}, ensuring a clear separation between the graphical user interface, the underlying sorting algorithms, and the application logic.

\subsection*{Structural Components}
\begin{enumerate}[label=\arabic*., itemsep=0.2em]
    \item \textbf{View (\texttt{sortamino-view.fxml} \& \texttt{style.css})}: Defines the user interface declaratively using FXML. The layout utilizes a \texttt{BorderPane} to organize controls: algorithm selection on the left, configuration on the top, action buttons and file I/O on the right, and the visualization/aggregation area in the center. A custom \texttt{style.css} handles all aesthetic styling, completely decoupling presentation from logic.
    \item \textbf{Controller (\texttt{Controller.java})}: Acts as the bridge between the View and the underlying services. It handles all FXML action events (button clicks, selections, window resizing) and delegates heavy lifting to utility classes to maintain a slim, readable class structure, adhering to the Single Responsibility Principle.
    \item \textbf{Models \& Utilities}:
    \begin{itemize}[itemsep=0.1em]
        \item \textbf{Algorithms Package}: Contains the \texttt{SortingStrategy} interface and its concrete implementations. The \texttt{StrategyType} enum acts as a registry.
        \item \textbf{Events Package}: Employs an event-sourcing pattern (\texttt{SwapEvent}, \texttt{CompareEvent}, etc.).
        \item \textbf{Array Generators}: Handles the creation of different initial array states using the Factory pattern.
        \item \textbf{SortingAnimator}: A static utility class responsible for generating JavaFX animations.
        \item \textbf{ComparisonService}: Handles the background execution of multiple algorithms and constructs the resulting data tables and charts.
        \item \textbf{FileService}: An isolated utility for handling file I/O operations.
    \end{itemize}
\end{enumerate}

\section{Design Decisions}

Several critical design decisions were made to prioritize performance, maintainability, and user experience:

\begin{itemize}[itemsep=0.3em]
    \item \textbf{Event-Sourcing Animation System}: Instead of attempting to pause algorithm execution using thread sleeps to wait for animations, the sorting algorithms execute at full CPU speed in the background and record a list of \texttt{SortingEvent} objects. The Controller then plays back these events using JavaFX transitions. This avoids UI thread blocking and allows dynamic speed adjustment.
    \item \textbf{Controller Decomposition}: The central \texttt{Controller.java} was refactored by extracting animation generation, comparison logic, and file handling into separate service classes. This prevents the "God Object" anti-pattern and makes the codebase highly modular.
    \item \textbf{Concurrency Management}: Both visualization and comparison modes execute heavily taxing logic on newly spawned background threads. UI updates are securely marshaled back to the JavaFX Application Thread using \texttt{Platform.runLater()}, preventing the application UI from freezing while processing large arrays.
    \item \textbf{Dual-Mode Interface}: Visualization Mode and Comparison Mode share the same FXML scene. The center pane dynamically swaps its children (animated rectangles vs. a data table) based on the user's selected mode, providing a seamless user experience.
\end{itemize}

\section{Sorting Comparison Output and Analysis}

The comparison mode evaluates all implemented algorithms across a user-defined array size and averages the results over a specified number of runs. For each run, a fresh array is generated to ensure fairness.

\subsection*{Expected Data Aggregation Results}

When analyzing a randomly generated array of size $N = 1000$ over several runs:

\begin{itemize}[itemsep=0.3em]
    \item \textbf{Selection Sort \& Bubble Sort}: Expected to show the highest average runtime. Comparisons will be exactly $\frac{N(N-1)}{2}$ for Selection Sort regardless of the input.
    \item \textbf{Insertion Sort}: While generally $\mathcal{O}(N^2)$ on random data, it drastically outperforms Bubble and Selection sorts on ``Almost Sorted'' arrays, approaching $\mathcal{O}(N)$ runtime and minimizing interchanges.
    \item \textbf{Merge Sort}: Consistently demonstrates low runtime due to its $\mathcal{O}(N \log N)$ complexity. The number of comparisons is steady, though array copies result in higher interchange counts compared to in-place sorts.
    \item \textbf{Quick Sort}: Generally the fastest algorithm for average runtime on random data. However, on ``Reversed'' arrays with naive pivoting, runtime and comparisons spike, illustrating its worst-case $\mathcal{O}(N^2)$ behavior.
    \item \textbf{Heap Sort}: Shows consistent $\mathcal{O}(N \log N)$ performance across all data types, proving its reliability as an in-place sort without the worst-case risks associated with Quick Sort.
\end{itemize}

\subsection*{Data Visualization}

The application generates two forms of data visualization upon concluding a comparison:
\begin{enumerate}[itemsep=0.1em]
    \item \textbf{Data Table}: Displays exact integers for Average Runtime (ms), Min/Max Runtime, Comparisons, and Interchanges.
    \item \textbf{Bar Chart}: Groups three data series side-by-side per algorithm: Average Runtime, Comparisons (scaled $\div 1000$), and Interchanges (scaled $\div 1000$). This visually highlights the stark contrast between $\mathcal{O}(N^2)$ and $\mathcal{O}(N \log N)$ scaling.
\end{enumerate}

\newpage
\section{Sample Outputs}

% INSTRUCTIONS FOR USER: 
% 1. Take the screenshots as described below.
% 2. Save them in the same folder as this .tex file.
% 3. Replace 'dashboard.png', 'sorting.png', etc. with your actual file names.
% 4. Uncomment the \includegraphics lines and recompile.

\subsection*{1. UI Dashboard (Idle State)}
% \begin{figure}[H]
%     \centering
%     \includegraphics[width=0.9\textwidth]{dashboard.png}
%     \caption{The application dashboard immediately after launch.}
% \end{figure}

\subsection*{2. Sorting Visualization (Active)}
% \begin{figure}[H]
%     \centering
%     \includegraphics[width=0.9\textwidth]{sorting.png}
%     \caption{Active sorting visualization showing color flashes for comparisons and geometry changes for swaps.}
% \end{figure}

\subsection*{3. File Input Integration}
% \begin{figure}[H]
%     \centering
%     \includegraphics[width=0.9\textwidth]{file_input.png}
%     \caption{UI state demonstrating a loaded custom array file with its element count.}
% \end{figure}

\subsection*{4. Comparison Mode Table}
% \begin{figure}[H]
%     \centering
%     \includegraphics[width=1\textwidth]{comparison_table.png}
%     \caption{Populated TableView showing aggregated metrics after completing runs.}
% \end{figure}

\subsection*{5. Comparison Chart}
% \begin{figure}[H]
%     \centering
%     \includegraphics[width=0.9\textwidth]{comparison_chart.png}
%     \caption{BarChart visualization highlighting algorithmic efficiency differences.}
% \end{figure}

\end{document}
